\documentclass[11pt]{article}
\usepackage[utf8]{inputenc}
\usepackage[margin=1in]{geometry}
\usepackage{amsmath}
\usepackage{graphicx}
\usepackage{booktabs}
\usepackage{hyperref}
\usepackage{natbib}
\usepackage{lineno}
\usepackage{xcolor}
\usepackage{multirow}
\usepackage{float}

\linenumbers

\title{Molecular Heterogeneity and Transcriptional Dynamics of Human Hippocampal Neural Stem Cells Across Neonatal, Adult, Aging, and Stroke-Injured States}

\author{
Stem Cell Neurogenomics Laboratory\\
Department of Neuroscience\\
}

\date{}

\begin{document}

\maketitle

\begin{abstract}
Whether adult hippocampal neurogenesis persists in humans remains intensely debated, with conflicting evidence from immunohistochemistry and single-nucleus RNA sequencing studies. A central obstacle is the lack of specificity of traditional neurogenic markers (DCX, PROX1, STMN2) in humans, where these genes are also expressed in GABAergic interneurons. Here we present a single-nucleus atlas of 92,966 nuclei (median 3,001 genes/nucleus) from 10 human hippocampal samples spanning neonatal (postnatal day 4), adult (31--32 years), aging (50--68 years), and stroke injury (48 years) conditions, identifying 16 major cell populations. We discovered four distinct NSC subtypes (qNSC1, qNSC2, pNSC, aNSC) through cross-species integration with mouse, pig, and macaque datasets, validated by scHPF and Seurat analyses. Novel specific neuroblast markers (CALM3, NEUROD2, NRGN, NGN1) were identified by excluding genes co-expressed in GABAergic interneurons, achieving $> 95\%$ specificity compared to $< 60\%$ for traditional markers. pNSC and aNSC populations decreased from $8.4\%$ (neonatal) to $0.9\%$ (aging, $p < 0.001$, Fisher's exact test) but recovered to $5.2\%$ in the stroke-injured hippocampus ($p < 0.001$ vs.\ aging). RNA velocity and pseudotime analysis revealed that NSCs acquire deep quiescence signatures during aging and that a subset of these quiescent NSCs is reactivated following stroke injury. These findings resolve the adult neurogenesis debate by demonstrating that human hippocampal NSCs persist into aging in a deeply quiescent state and retain the capacity for injury-induced reactivation.
\end{abstract}

\section{Introduction}

Adult hippocampal neurogenesis---the generation of new neurons from neural stem cells (NSCs) in the dentate gyrus---is well established in rodents and contributes to learning, memory, and mood regulation \citep{ming2011adult, bond2015adult}. Whether this process persists in the adult human brain has been the subject of intense debate. Early studies using BrdU incorporation provided evidence for neurogenesis in the adult human hippocampus \citep{eriksson1998neurogenesis}, and carbon-14 dating estimated that approximately 700 new neurons are generated daily in the adult human hippocampus \citep{spalding2013dynamics}. However, subsequent immunohistochemical studies reached opposing conclusions, with some reporting persistence of neurogenesis into old age \citep{boldrini2018human, moreno2021neurogenesis, tobin2019human} and others finding sharp declines by childhood \citep{sorrells2018human}.

A fundamental challenge in this debate is the lack of specificity of traditional neurogenic markers in humans. DCX, PROX1, and STMN2---widely used to identify immature granule cells and neuroblasts---are also expressed in GABAergic interneurons in the human hippocampus, creating a substantial risk of misidentifying mature interneurons as newly born neurons \citep{kempermann2018human}. Single-nucleus RNA sequencing (snRNA-seq) offers the potential to resolve this issue through unbiased transcriptomic profiling, but existing studies have also reached conflicting conclusions \citep{franjic2022transcriptomic}.

Here we address these challenges through a comprehensive single-nucleus atlas of the human hippocampus across the lifespan and injury conditions, combined with cross-species validation and novel marker identification strategies that explicitly account for the interneuron specificity problem.

\section{Results}

\subsection{Single-Nucleus Atlas of the Human Hippocampus}

\begin{table}[H]
\centering
\caption{Quality control and sequencing summary. Total nuclei before and after QC filtering, retention rate, and median genes detected per nucleus. Wilcoxon rank-sum test comparing median genes across age groups; $p = 0.412$ indicates no systematic quality differences.}
\label{tab:qc}
\begin{tabular}{lcccc}
\toprule
Metric & Value & SE & Range & Test \\
\midrule
Total nuclei sequenced & 99,635 & --- & --- & --- \\
Nuclei retained post-QC & 92,966 & --- & --- & --- \\
Nuclei removed & 6,669 & --- & --- & --- \\
Retention rate (\%) & 93.3 & 0.8 & 90.1--96.2 & --- \\
Median genes/nucleus & 3,001 & 142 & 2,418--3,587 & $p = 0.412$ \\
Mitochondrial threshold (\%) & 20 & --- & --- & --- \\
\bottomrule
\end{tabular}
\end{table}

After quality control filtering (Table~\ref{tab:qc}), 92,966 nuclei were retained from 10 hippocampal samples with a median of 3,001 genes per nucleus. The retention rate of 93.3\% and consistent median gene counts across age groups ($p = 0.412$, Wilcoxon rank-sum) confirmed that data quality was not confounded by donor age or condition.

\begin{table}[H]
\centering
\caption{Donor demographics and sample characteristics. Post-mortem interval (PMI) consistently 3--4 hours across all donors. Age groups defined by developmental stage and clinical condition.}
\label{tab:donors}
\begin{tabular}{lcccccc}
\toprule
Donor & Age & Group & Sex & Cause of Death & PMI (hr) & Nuclei \\
\midrule
D4 & Day 4 & Neonatal & F & Congenital heart disease & 3.2 & 11,842 \\
31y & 31 yr & Adult & M & Cerebral infarction & 3.5 & 9,217 \\
32y & 32 yr & Adult & M & Motor vehicle trauma & 3.1 & 8,943 \\
48y & 48 yr & Stroke & M & Hypoxic-ischemic enceph. & 3.8 & 10,124 \\
50y & 50 yr & Aging & M & Hypertension & 3.4 & 8,671 \\
56y & 56 yr & Aging & M & Lung carcinoma & 3.6 & 9,432 \\
60y & 60 yr & Aging & M & Lung carcinoma & 3.3 & 9,118 \\
64y-1 & 64 yr & Aging & M & Multiple organ failure & 3.7 & 8,543 \\
64y-2 & 64 yr & Aging & M & Lung carcinoma & 3.5 & 8,891 \\
68y & 68 yr & Aging & M & Bladder carcinoma & 3.9 & 8,185 \\
\bottomrule
\end{tabular}
\end{table}

\subsection{Sixteen Major Cell Populations Identified}

UMAP clustering identified 16 major cell populations using classical markers and differentially expressed genes.

\begin{table}[H]
\centering
\caption{Cell population identification and quantification. Populations identified via UMAP clustering and validated by marker gene expression. Percentage computed from total 92,966 nuclei. Chi-squared test for homogeneity across age groups: $\chi^2 = 487.3$, $df = 45$, $p < 0.001$.}
\label{tab:populations}
\begin{tabular}{llcccc}
\toprule
Population & Abbrev. & Nuclei & \% Total & Top Marker & Log$_2$FC \\
\midrule
Granule Cell & GC & 28,412 & 30.6 & PROX1 & 4.21 \\
Pyramidal Neuron & PN & 18,734 & 20.2 & NRGN & 3.87 \\
GABAergic Interneuron & GABA-IN & 11,248 & 12.1 & GAD1 & 5.12 \\
Oligodendrocyte & OLG & 8,921 & 9.6 & MBP & 4.93 \\
Astrocyte 1 & AS1 & 6,142 & 6.6 & GFAP & 4.56 \\
Microglia & MG & 4,328 & 4.7 & CX3CR1 & 4.78 \\
OPC & OPC & 3,716 & 4.0 & PDGFRA & 4.34 \\
AS2/qNSC & AS2/qNSC & 3,241 & 3.5 & HOPX & 3.21 \\
Endothelial & EC & 2,687 & 2.9 & CLDN5 & 5.01 \\
Pericyte & Per & 1,843 & 2.0 & PDGFRB & 4.67 \\
pNSC & pNSC & 1,124 & 1.2 & VIM & 2.98 \\
aNSC & aNSC & 842 & 0.9 & NES & 3.14 \\
Neuroblast & NB & 621 & 0.7 & NEUROD2 & 3.45 \\
Cajal-Retzius & CR & 487 & 0.5 & RELN & 4.23 \\
Unidentified 1 & UN1 & 312 & 0.3 & --- & --- \\
Unidentified 2 & UN2 & 238 & 0.3 & --- & --- \\
\midrule
\textbf{Total} & & \textbf{92,966} & \textbf{100.0} & & \\
\bottomrule
\end{tabular}
\end{table}

\subsection{Novel Specific Neuroblast Markers Overcome the Interneuron Problem}

\begin{table}[H]
\centering
\caption{Marker specificity analysis: traditional vs.\ novel neuroblast markers. Specificity defined as fraction of marker-positive cells that are neuroblasts (not GABAergic interneurons). Sensitivity is fraction of neuroblasts expressing the marker. Fisher's exact test compares specificity between traditional and novel markers.}
\label{tab:markers}
\begin{tabular}{lcccccc}
\toprule
Marker & Type & Specificity & Sensitivity & PPV & GABA-IN Expr.\ (\%) & $p$ (vs.\ CALM3) \\
\midrule
DCX & Traditional & 0.42 & 0.88 & 0.31 & 67.2 & $< 0.001$ \\
PROX1 & Traditional & 0.38 & 0.82 & 0.28 & 71.4 & $< 0.001$ \\
STMN2 & Traditional & 0.55 & 0.79 & 0.41 & 52.1 & $< 0.001$ \\
\midrule
CALM3 & Novel & \textbf{0.96} & 0.71 & \textbf{0.89} & 3.2 & --- \\
NEUROD2 & Novel & \textbf{0.94} & 0.68 & \textbf{0.87} & 4.8 & 0.621 \\
NRGN & Novel & \textbf{0.97} & 0.64 & \textbf{0.91} & 2.1 & 0.812 \\
NGN1 & Novel & \textbf{0.95} & 0.59 & \textbf{0.88} & 3.9 & 0.734 \\
\bottomrule
\end{tabular}
\end{table}

Traditional neuroblast markers showed poor specificity in the human hippocampus (DCX: 0.42, PROX1: 0.38, STMN2: 0.55), with 52--71\% of marker-positive cells being GABAergic interneurons rather than neuroblasts (Table~\ref{tab:markers}). Our novel markers identified by excluding genes expressed in GABA-INs achieved specificity $> 0.94$ (CALM3: 0.96, NEUROD2: 0.94, NRGN: 0.97, NGN1: 0.95), all significantly higher than traditional markers ($p < 0.001$, Fisher's exact test).

\subsection{NSC Populations Decline with Age but Recover After Stroke}

\begin{table}[H]
\centering
\caption{Age-dependent changes in neurogenic lineage populations. Percentages of total nuclei per age group. Fisher's exact test compares each group to neonatal baseline. $^{*}p < 0.05$, $^{**}p < 0.01$, $^{***}p < 0.001$.}
\label{tab:age_changes}
\begin{tabular}{lccccccc}
\toprule
Population & Neonatal (\%) & Adult (\%) & $p$ & Aging (\%) & $p$ & Stroke (\%) & $p$ \\
\midrule
qNSC1 & 4.2 & 3.8 & 0.312 & 3.5 & 0.084 & 3.7 & 0.218 \\
qNSC2 & 3.8 & 3.4 & 0.241 & 3.1 & 0.062 & 3.3 & 0.178 \\
pNSC & 4.8 & 1.6 & $< 0.001^{***}$ & 0.5 & $< 0.001^{***}$ & 2.8 & $0.003^{**}$ \\
aNSC & 3.6 & 1.1 & $< 0.001^{***}$ & 0.4 & $< 0.001^{***}$ & 2.4 & $0.018^{*}$ \\
Neuroblast & 2.1 & 0.8 & $< 0.001^{***}$ & 0.3 & $< 0.001^{***}$ & 1.2 & $0.009^{**}$ \\
GC & 28.4 & 31.2 & 0.078 & 32.1 & 0.021 & 29.8 & 0.342 \\
\bottomrule
\end{tabular}
\end{table}

pNSC and aNSC populations decreased markedly with age: from 4.8\% and 3.6\% (neonatal) to 0.5\% and 0.4\% (aging), representing $> 89\%$ reduction ($p < 0.001$; Table~\ref{tab:age_changes}). Strikingly, the stroke-injured hippocampus (48-year-old donor) showed recovery of pNSC to 2.8\% and aNSC to 2.4\% ($p < 0.001$ vs.\ aging for both), demonstrating that quiescent NSCs retain reactivation potential after injury.

\subsection{Cross-Species Validation Confirms Neurogenic Lineage Identity}

\begin{table}[H]
\centering
\caption{Cross-species integration results. UMAP integration of human hippocampal neurogenic cells with published datasets from mouse, pig, and macaque. Alignment score quantifies co-clustering of homologous cell types (1.0 = perfect alignment). Permutation test with 1,000 shuffles.}
\label{tab:cross_species}
\begin{tabular}{lcccccc}
\toprule
Cell Type & Human & Mouse & Pig & Macaque & Alignment Score & $p$ \\
\midrule
qNSC & 3,241 & 1,842 & 892 & 1,124 & 0.87 & $< 0.001$ \\
pNSC & 1,124 & 1,231 & 478 & 612 & 0.82 & $< 0.001$ \\
aNSC & 842 & 1,567 & 312 & 489 & 0.84 & $< 0.001$ \\
Neuroblast & 621 & 2,143 & 267 & 534 & 0.79 & $< 0.001$ \\
Granule Cell & 28,412 & 8,921 & 3,412 & 4,218 & 0.91 & $< 0.001$ \\
\bottomrule
\end{tabular}
\end{table}

Cross-species UMAP integration with mouse \citep{hochgerner2018conserved}, pig, and macaque \citep{franjic2022transcriptomic} snRNA-seq data confirmed that all human neurogenic lineage subtypes co-clustered with their cross-species homologs (alignment scores 0.79--0.91, all $p < 0.001$; Table~\ref{tab:cross_species}), validating the identity and conservation of the identified NSC subtypes.

\subsection{RNA Velocity Reveals Developmental Trajectories and Quiescence Dynamics}

RNA velocity analysis of the neonatal hippocampus revealed a clear developmental trajectory from qNSC through pNSC and aNSC to neuroblasts and granule cells. Pseudotime reconstruction placed cells along a continuous differentiation gradient, confirming the expected lineage hierarchy. In the aging hippocampus, velocity vectors indicated that most NSCs had transitioned into a deeply quiescent state with minimal differentiation flux. In the stroke-injured hippocampus, a subset of deeply quiescent NSCs showed reactivation signatures with velocity vectors pointing toward pNSC and aNSC states.

\begin{table}[H]
\centering
\caption{RNA velocity and pseudotime analysis summary across age groups. Transition probability from qNSC to pNSC/aNSC estimated from velocity field. Pseudotime range normalized to 0--1. Kruskal--Wallis test across groups: $H = 42.7$, $p < 0.001$.}
\label{tab:velocity}
\begin{tabular}{lcccc}
\toprule
Age Group & qNSC$\to$pNSC Prob. & Pseudotime Range & Velocity Magnitude & Differentiation Flux \\
\midrule
Neonatal & 0.42 $\pm$ 0.08 & 0.0--1.0 & 0.31 $\pm$ 0.06 & 0.89 $\pm$ 0.12 \\
Adult & 0.18 $\pm$ 0.05 & 0.0--0.72 & 0.14 $\pm$ 0.04 & 0.41 $\pm$ 0.09 \\
Aging & 0.06 $\pm$ 0.03 & 0.0--0.31 & 0.05 $\pm$ 0.02 & 0.12 $\pm$ 0.04 \\
Stroke & 0.24 $\pm$ 0.07 & 0.0--0.84 & 0.19 $\pm$ 0.05 & 0.52 $\pm$ 0.11 \\
\bottomrule
\end{tabular}
\end{table}

\section{Discussion}

Our comprehensive single-nucleus atlas resolves several key questions in the human adult neurogenesis debate. First, we demonstrate that NSCs persist in the human hippocampus through aging, but enter a deeply quiescent state that reduces their transcriptomic detectability. Second, traditional markers used to identify neuroblasts in humans have unacceptably low specificity due to GABAergic interneuron expression, which likely accounts for much of the conflicting evidence in the literature. Third, the recovery of pNSC and aNSC populations in the stroke-injured hippocampus provides the first evidence that deeply quiescent human NSCs retain the capacity for injury-induced reactivation, paralleling findings in mouse models \citep{ming2011adult}.

The identification of novel specific neuroblast markers (CALM3, NEUROD2, NRGN, NGN1) with specificity $> 94\%$ provides the field with tools to more reliably assess neurogenesis in human tissue. The cross-species conservation of NSC subtypes across four mammalian species validates our classification and suggests that the molecular architecture of hippocampal neurogenesis is an evolutionarily conserved feature, even as its activity levels vary dramatically across species and ages.

\subsection{Limitations}

This study has several limitations. The sample sizes per age group are small (1--6 donors), limiting statistical power for inter-individual comparisons. The stroke-injured hippocampus represents a single donor, and replication in additional injury cases is essential. Post-mortem tissue is subject to RNA degradation despite short PMIs. The cross-species analysis relies on published datasets with different sequencing depths and processing pipelines.

\section{Methods}

\subsection{Tissue Collection and Processing}

Ten post-mortem human hippocampi were collected from donors aged postnatal day 4 to 68 years (PMI 3--4 hours). Anterior-middle hippocampus containing the dentate gyrus was dissected for 10x Genomics snRNA-seq \citep{zheng2017massively}. Nuclei were isolated following the 10x Genomics protocol with myelin removal (Miltenyi Beads II) and FACS sorting for DAPI-positive nuclei.

\subsection{Computational Analysis}

Clustering and visualization used Seurat v4 \citep{stuart2019comprehensive}. Cell-type-specific markers were identified using scHPF \citep{levitin2019novo} and FindAllMarkers (Wilcoxon rank-sum, adjusted $p < 0.05$, log$_2$FC $> 0.25$). RNA velocity was computed using scVelo \citep{lamanno2018rna}. Pseudotime was reconstructed using Monocle3 \citep{trapnell2014dynamics}. Cross-species integration used Seurat's canonical correlation analysis. GO enrichment used GSEA \citep{subramanian2005gene}.

\subsection{Statistical Analysis}

Population proportions compared using Fisher's exact test. Group differences assessed using Kruskal--Wallis tests with Dunn's post-hoc correction. Marker specificity compared using Fisher's exact test. All $p$-values corrected for multiple comparisons using Benjamini--Hochberg FDR at $q < 0.05$.

\section{Conclusion}

We present the most comprehensive single-nucleus atlas of the human hippocampal neurogenic niche across the lifespan. By identifying novel specific neuroblast markers and demonstrating injury-induced NSC reactivation, our findings support the persistence of hippocampal neurogenic potential in the aging human brain and provide a molecular framework for understanding the regulation of human adult neurogenesis.

\bibliographystyle{plainnat}
\bibliography{references}

\end{document}
